\chapter{Proposed System Methodology (Sentinel Forge)}

\section{Conceptual Framework and Design Philosophies}

Sentinel Forge is built on three core design philosophies:

\begin{enumerate}
    \item \textbf{Detect Out-of-Band, Respond In-Band:} The sensor analyzes copies of traffic asynchronously; enforcement actions are executed on the target infrastructure through authenticated remote commands.
    
    \item \textbf{Defense in Depth:} No single detection layer is trusted. Packets pass through fragment reassembly, multi-format decoding, entropy analysis, ML triage, and weighted rule matching before classification.
    
    \item \textbf{Zero-Trust Inter-Component Communication:} Every command between the sensor and the enforcement agent must be cryptographically verified, preventing unauthorized blocking actions.
\end{enumerate}

% ---- FIGURE: High-Level System Architecture ----
\begin{figure}[h]
\centering
\fbox{\parbox{0.85\textwidth}{\centering\vspace{4cm}\textit{Insert Diagram: High-level Sentinel Forge system architecture block diagram showing all 6 modules --- Static Analyzer, Malware Simulator, Dynamic Analyzer (DPI Engine), Inference Engine, Universal Agent Core, and Threat Intel Module --- with arrows showing data flow between them}\vspace{4cm}}}
\caption{Sentinel Forge: High-Level System Architecture}
\label{fig:system_architecture}
\end{figure}

\section{Network Topology: The SPAN/Mirror Port Architecture}

Sentinel Forge deploys as an OOB sensor connected to a SPAN (Switched Port Analyzer) or mirror port on the network switch. The sensor receives an exact copy of all traffic traversing the monitored segment. A separate channel connects the sensor to the Universal Agent Core residing on (or accessible to) the target infrastructure.

% ---- FIGURE: Network Topology ----
\begin{figure}[h]
\centering
\fbox{\parbox{0.85\textwidth}{\centering\vspace{3cm}\textit{Insert Screenshot: VirtualBox Lab Network Topology --- showing Attacker VM, pfSense Firewall VM, Target Server VM, and Sentinel Forge Sensor VM with IP addresses and network interfaces labeled}\vspace{3cm}}}
\caption{Sentinel Forge Network Topology with SPAN Port Architecture}
\label{fig:network_topology}
\end{figure}

\section{The Sentinel Forge Pipeline}

% ---- FIGURE: DPI Pipeline Flowchart ----
\begin{figure}[h]
\centering
\fbox{\parbox{0.85\textwidth}{\centering\vspace{4cm}\textit{Insert Flowchart: The complete DPI pipeline --- Packet Ingestion $\rightarrow$ Fragment Reassembly $\rightarrow$ Layer Re-dissection $\rightarrow$ Normalization (Raw/B64/Hex/XOR) $\rightarrow$ ML Pre-Filter (Safe? $\rightarrow$ Benign Buffer; Suspicious? $\rightarrow$ continue) $\rightarrow$ Rule Matching $\rightarrow$ Score $\geq$ Threshold? $\rightarrow$ Alert + Quarantine + Webhook Dispatch}\vspace{4cm}}}
\caption{Sentinel Forge Detection Pipeline Flowchart}
\label{fig:pipeline_flowchart}
\end{figure}

The detection pipeline processes each packet through the following stages:

\subsection{Traffic Ingestion and IP Defragmentation}

The \texttt{FragmentReassembler} module buffers incoming IP fragments indexed by the tuple (source IP, destination IP, IP ID). When the final fragment arrives (MF flag = 0), Scapy's \texttt{defrag()} function reconstructs the complete packet. After reassembly, the TCP/UDP layer is forcefully re-parsed from raw IP bytes to recover transport-layer fields that Scapy may have missed during fragmentation.

\subsection{Multi-Layer Decryption and Normalization}

The Packet Normalizer generates multiple representations of each payload for anti-evasion matching:

\begin{enumerate}
    \item \textbf{Raw:} Original payload bytes.
    \item \textbf{Base64 Decoded:} Attempts standard Base64 decoding.
    \item \textbf{Hex Decoded:} Attempts hexadecimal string-to-bytes conversion.
    \item \textbf{XOR Decoded:} Brute-forces all 256 single-byte XOR keys. Only the first key producing $>$30\% printable ASCII characters is retained, preventing storage of 255 garbage variants.
\end{enumerate}

Each decoded variant is checked during rule matching, ensuring detection regardless of the obfuscation layer applied by the attacker.

% ---- FIGURE: Anti-Evasion Layers ----
\begin{figure}[h]
\centering
\fbox{\parbox{0.85\textwidth}{\centering\vspace{3cm}\textit{Insert Diagram: Anti-evasion pipeline layers --- showing a single incoming packet being decoded into 4 variants (Raw, Base64, Hex, XOR) like a funnel, with each variant checked against rules}\vspace{3cm}}}
\caption{Multi-Layer Anti-Evasion Payload Decoding Pipeline}
\label{fig:antievasion_pipeline}
\end{figure}

\subsection{Shannon Entropy Analysis for Anomaly Detection}

Shannon Entropy quantifies the randomness of a byte sequence \cite{shannon1948}:

\begin{equation}
H(X) = -\sum_{i=0}^{255} p(x_i) \cdot \log_2 p(x_i)
\end{equation}

Values range from 0.0 (uniform) to 8.0 (perfectly random). Encrypted or compressed payloads exhibit entropy $>$ 7.5, making this metric effective for detecting obfuscated rootkit communications. Entropy is calculated only on payloads $\geq$ 32 bytes to ensure statistical significance. High entropy triggers a 35-point boost in the risk scoring engine.

\section{AI-Driven Triage: The Decision Tree Inference Model}

Before full DPI analysis, each packet is evaluated by a Decision Tree classifier (max\_depth = 5) trained on five features:

\begin{table}[h]
\centering
\caption{ML Pre-Filter Feature Set}
\label{tab:ml_features}
\begin{tabular}{|c|l|l|}
\hline
\textbf{No.} & \textbf{Feature} & \textbf{Type} \\
\hline
1 & Shannon Entropy & Continuous (0.0 -- 8.0) \\
\hline
2 & Destination Port & Integer \\
\hline
3 & Source Port & Integer \\
\hline
4 & TCP Window Size & Integer \\
\hline
5 & Protocol (TCP=1, Other=0) & Binary \\
\hline
\end{tabular}
\end{table}

Packets classified as safe (prediction = 0) bypass the rule engine entirely and are saved to the benign buffer for continuous baseline updating. This enables the system to discard the majority of benign traffic in sub-millisecond time.

% ---- FIGURE: Decision Tree Visualization ----
\begin{figure}[h]
\centering
\fbox{\parbox{0.85\textwidth}{\centering\vspace{3cm}\textit{Insert Diagram: Visualization of the trained Decision Tree (max\_depth=5) showing the decision splits on Entropy, Port, Window Size --- can be generated using sklearn's export\_graphviz or plot\_tree}\vspace{3cm}}}
\caption{Decision Tree Pre-Filter: Trained Model Visualization}
\label{fig:decision_tree_viz}
\end{figure}

\section{Secure Response Orchestration: HMAC-Signed Webhook Protocols}

When a threat is detected, the sensor constructs a JSON payload containing the target IP, action (``block''), and timestamp, then signs it using HMAC-SHA256 with a pre-shared secret key. The signed webhook is dispatched asynchronously via a daemon thread to the Universal Agent Core.

The Agent Core validates the signature using a timing-attack-resistant comparison (\texttt{hmac.compare\_digest}), checks the IP against a configurable whitelist, queries the idempotency cache (SQLite), and delegates enforcement to the loaded firewall plugin. This multi-layered validation prevents unauthorized blocking, duplicate actions, and accidental self-disruption.

% ---- FIGURE: HMAC Webhook Sequence ----
\begin{figure}[h]
\centering
\fbox{\parbox{0.85\textwidth}{\centering\vspace{3cm}\textit{Insert Diagram: Sequence diagram showing: Sensor detects threat $\rightarrow$ Signs payload with HMAC-SHA256 $\rightarrow$ Sends POST /api/v1/alert $\rightarrow$ Agent Core verifies signature $\rightarrow$ Checks whitelist $\rightarrow$ Checks idempotency cache $\rightarrow$ Calls firewall plugin $\rightarrow$ Returns result}\vspace{3cm}}}
\caption{HMAC-SHA256 Webhook Authentication Sequence}
\label{fig:hmac_sequence}
\end{figure}

An LRU rate-limiting cache with configurable TTL prevents webhook storms during volumetric attacks by enforcing a cooldown period per source IP.
